\section{2018秋}

\setcounter{yearcounter}{2018}

\subsection{微分積分}
以下の問いに答えなさい。
\prob{
  \vskip.5\baselineskip
  点$(5,\,0)$から曲線$2y^2 = x^3$まで最短距離を求めよ。
  \vskip.5\baselineskip
}
\begin{ans*}
  ${}$
  距離$L$は$2y^2 = x^3$上の点$(x,\,y)$に依存するので$L = L(x,y)$と表せる。
  このとき
  \begin{align}
    L^2
    &= (x - 5)^2 + y^2 \\
    &= \frac{x^3}{2} + x^2 - 10x + 25
  \end{align}
  であって、定義域は$2y^2 = x^3$より$x\geq 0$である。
  この最小値を求める。
  \begin{align}
    \frac{d}{dx}L^2
    = \frac{1}{2}(3x^2 + 4x - 20)
    = \frac{1}{2}(x - 2)(3x + 10)
  \end{align}
  より、\dm{\frac{d(L^2)}{dx} = 0\land x\geq 0\Rightarrow x = 2}で
  $L^2$の増減は下図。
  \begin{table}[H]
  \centering
  \begin{tabular}{c||cccc}
  \toprule
    $x$ & 0 & $\cdots$ & 2 & $\cdots$ \\
  \midrule
    \dm{\frac{d(L^2)}{dx}} & & $-$ & 0 & $+$ \\[.8em]
    \dm{L^2} & 25 & \searrow & 13 & \nearrow \\[.3em]
    \dm{L} &&& $\sqrt{13}$ & \\
  \bottomrule
  \end{tabular}\end{table}
  ゆえに、求める最短距離は$\sqrt{13}$。
\end{ans*}

\prob{
  次の積分の値を各領域$D$において求めなさい。
  \begin{enumerate}[label=(\arabic*), ref=(\arabic*), itemsep=0pt]
    \item \dm{D = \Dset{(x,y)\relmiddle 1\leq y\leq 2,\,y\leq x\leq y^3}}
    \begin{gather}
      \iint_{D} e^{\frac{x}{y}}\,dxdy
    \end{gather}
    \item 領域\dm{D}は$y = \sqrt{x}$と$y = x^3$に囲まれた領域とせよ。
    \begin{gather}
      \iint_{D}(4xy - y^3)dxdy
    \end{gather}
  \end{enumerate}
}
\begin{ans*}
  ${}$
  \begin{enumerate}[label=(\arabic*), ref=(\arabic*), itemsep=0pt]
    \item ${}$
    \begin{align}
      \iint_D e^{\frac{x}{y}}\,dxdy
      &= \int_{1}^{2}dy \biggl[ye^\frac{x}{y}\biggr]_{y}^{y^3} \\
      &= \int_{1}^{2}dy \Bigl(ye^{y^2} - ey\Bigr) \\
      &= \int_{1}^{2}ye^{y^2} \,dy - e\int_{1}^{2}y\,dy \\
      &= \frac{1}{2}\biggl[e^{y^2}\biggr]_{1}^{2} - \frac{e}{2}\biggl[y^{2}\biggr]_{1}^{2} \\
      &= \frac{1}{2}(e^{4} - e) - \frac{3}{2}e \\
      &= \frac{1}{2}e^{4} - 2e
    \end{align}
    \item ${}$
    \begin{align}
      \iint_D (4xy - y^3)dxdy
      &= \int_{0}^{1}\int_{x^3}^{\sqrt{x}}(4xy - y^3)dydx \\
      &= \int_{0}^{1}dx \biggl[2xy^2 - \frac{1}{4}y^4\biggr]_{x^3}^{\sqrt{x}} \\
      &= \int_{0}^{1}\Bigl(2x^2 - \frac{x^2}{4} - 2x^7 + \frac{x^{12}}{4}\Bigr) \\
      &= \biggl[\frac{7}{12}x^3 - \frac{1}{4}x^8 + \frac{1}{56}x^{13}\biggr]_{0}^{1} \\
      &= \frac{55}{156}
    \end{align}
  \end{enumerate}
\end{ans*}

\prob{
  次の微分方程式を解きなさい。
  ただし、$x(0) = x_{0}$および$\dot{x}(0) = u_{0}$とせよ。
  \begin{gather}
    \ddot{x} - x = 0
  \end{gather}
  ここで、\dm{\dot{x} = \frac{dx}{dt}}および
  \dm{\ddot{x} = \frac{d^2 x}{dt^2}}である。
}
\begin{ans*}
  補助方程式$\grl^2 - 1 = 0\Leftrightarrow \grl = \pm 1$より、
  一般解は$A,\,B$を任意定数として\dm{y = Ae^{t} + Be^{-t}}と表せる。

  ここに、初期条件
  \begin{gather}
    A + B = x_0,\qquad
    A - B = u_0
  \end{gather}
  を考えれば、初期条件を考慮した一般解は
  \begin{gather}
    y = \frac{x_0 + u_0}{2}e^{t} + \frac{x_0 - u_0}{2}e^{-t}
  \end{gather}
\end{ans*}

\newpage
\subsection{線形代数}

\prob{
  次の行列$\bA$について、以下の問いに答えなさい。
  \begin{gather}
    \bA = \bmat{
      10 & 4 & 0 \\
      4 & -5 & 0 \\
      0 & 0 & \beta
    }
  \end{gather}
  ここで、$\beta$は実数である。
  また、固有値の重複はないものとする。
  \begin{enumerate}[label=(\arabic*), ref=(\arabic*), itemsep=0pt]
    \item 固有値を求めなさい。
    \item 正規化された固有ベクトルを求めなさい。
    \item $\bA$が正則のときの条件を答え、逆行列を求めなさい。
    \item $\bA$が正則でないとき、階数を求めなさい。
    \item \label{item:2018a-l-1-5}$\bA$の部分行列
    \begin{gather}
      \bB = \bmat{
        10 & 4 \\
        4 & 5
      }
    \end{gather}
    について、$\bB^2 + p\bB + q\bI = \bzv$が成り立つとき、
    実数$p$と$q$を求めなさい。
    ただし、$\bzv$は零行列、$\bI$は単位行列である。
    \item \ref*{item:2018a-l-1-5}で与えた$\bA$の部分行列$\bB$について、
    $\bB^n$の1行1列の成分を求めなさい。
    ただし、$n$は自然数である。
  \end{enumerate}
}

\begin{ans*}
  ${}$ \\
  \begin{enumerate}[label=(\arabic*), ref=(\arabic*), itemsep=0pt]
    \item 固有方程式\dm{(10-\grl)(-5-\grl)(\beta-\grl) - 16(\beta-\grl)=0}より、
    固有値$\grl$は
    \begin{gather}
      (10-\grl)(-5-\grl)(\beta-\grl) - 16(\beta-\grl)
      = -(\grl - \beta)(\grl - 11)(\grl + 6) = 0 \\
      \therefore \grl = \beta,\,11,\,-6
    \end{gather}
    \item 固有ベクトルを$\bu$とする。
    \begin{enumerate}[label=(\roman*), ref=(\roman*), itemsep=0pt]
      \item $\grl=\beta$のとき、
      \begin{gather}
        (\bA - \grl\bE)\bu = \bmat{
          10-\beta & 4 & 0 \\
          4 & -5-\beta & 0 \\
          0 & 0 & 0
        }\bu = \bzv \\
        \therefore\bu_1 = \Bmat{
          0 \\ 0 \\ 1
        }
      \end{gather}
      \item $\grl=11$のとき、
      \begin{gather}
        (\bA - \grl\bE)\bu = \bmat{
          -1 & 4 & 0 \\
          4 & -16 & 0 \\
          0 & 0 & \beta-11
        }\bu = \bzv \\
        \therefore\bu_2 = \frac{1}{\sqrt{17}}\Bmat{
          -1 \\ 4 \\ 0
        }
      \end{gather}
      \item $\grl=-6$のとき、
      \begin{gather}
        (\bA - \grl\bE)\bu = \bmat{
          16 & 4 & 0 \\
          4 & 1 & 0 \\
          0 & 0 & \beta+6
        }\bu = \bzv \\
        \therefore\bu_3 = \frac{1}{\sqrt{17}}\Bmat{
          -1 \\ 4 \\ 0
        }
      \end{gather}
    \end{enumerate}
    \item ${}$
    \begin{align}
      \bA\text{が正則}
      &\Leftrightarrow \det\bA \neq 0 \\
      &\Leftrightarrow -50\beta-16\beta \neq 0 \\
      &\Leftrightarrow \beta \neq 0
    \end{align}

    また、$\beta\neq 0$のとき、
    \begin{align}
      \left[\begin{array}[]{ccc|ccc}
        10 & 4 & 0 & 1 & 0 & 0 \\
        4 & -5 & 0 & 0 & 1 & 0 \\
        0 & 0 & \beta & 0 & 0 & 1
      \end{array}\right]
      &\longrightarrow
      \left[\begin{array}[]{ccc|ccc}
        10 & 4 & 0 & 1 & 0 & 0 \\
        0 & 1 & 0 & 4/66 & -10/66 & 0 \\
        0 & 0 & 1 & 0 & 0 & 1/\beta
      \end{array}\right]\\
      &\longrightarrow
      \left[\begin{array}[]{ccc|ccc}
        &&& 5/66 & 4/66 &  0 \\
        &\bE&& 4/66 & -10/66 & 0 \\
        &&& 0 & 0 & 1/\beta
      \end{array}\right]
    \end{align}
    ゆえに、$\bA$の逆行列は
    \begin{gather}
      \bA^{-1} = \bmat{
        5/66 & 2/33 &  0 \\
        2/33 & -5/33 & 0 \\
        0 & 0 & 1/\beta
      }
    \end{gather}
    \item $\beta = 0$のとき、
    \begin{align}
      \bA
      \longrightarrow\bmat{
        10 & 4 & 0 \\
        0 & -33/5 & 0 \\
        0 & 0 & 0
      }
    \end{align}
    より、$\rank\bA = 2$
    \item この部分行列$\bB$の固有方程式は\dm{(10-\grl)(-5-\grl) - 16 = \grl^{2} - 5\grl - 66 = 0}
    であるので、ケーリーハミルトンの定理から % [?] これ以外にないことは言えてる？
    \begin{gather}
      \bB^{2} -5\bB - 66\bI = \bzv
    \end{gather}
    ゆえに、
    \begin{gather}
      p = -5 ,~ q = -66
    \end{gather}
    \item $\bB^{n}$は$\bB$についての多項式$Q(\bB)$を用いて、
    \begin{gather}
      \bB^{n} = Q(\bB)(\bB^2 - 5\bB - 66) a\bB + b\bI \label{eq:2018a_l_1}
    \end{gather}
    と表される。ここに、$x\in\bbR$を用いて
    \begin{gather}
      x^{n} = Q(x)(x^2 - 5x - 66) + ax + b
    \end{gather}
    と書けるが、$x = 11,\,-6$として
    \begin{gather}
      a = \frac{11^{n} - (-6)^{n}}{17} \\
      b = \frac{6\cdot 11^{n} + 11\cdot(-6)}{17}
    \end{gather}
    ゆえに、\refeq{eq:2018autumn_l_1}の$a,\,b$はこのように与えられて、
    $\bB^n$の1行1列の成分$B^{n}_{\:11}$は
    \begin{align}
      B^{n}_{\:11}
      &= aB_{11} + b \\
      &= \frac{1}{17} \bigl(16\cdot 11^{n} + (-6)^{n}\bigr)
    \end{align}
    \begin{other*}
      $\bB$は$\grl = -6,\,11$の固有値をもつ。固有ベクトルを$\bp$で表すとき、
      \begin{enumerate}[label=(\roman*), ref=(\roman*), itemsep=0pt]
        \item $\grl = -6$のとき、
        \begin{gather}
          \bmat{
            16 & 4 \\
            4 & 1
          }\bp = \bzv \\
          \therefore \bp_1 = \Bmat{
            1 \\ -4
          }
        \end{gather}
        \item $\grl = 11$のとき、
        \begin{gather}
          \bmat{
            -1 & 4 \\
            4 & -16
          }\bp = \bzv \\
          \therefore \bp_2 = \Bmat{
            4 \\ 1
          }
        \end{gather}
      \end{enumerate}
      これらの固有ベクトルを並べた新たな行列$\grL:=[\bp_1,\,\bp_2]$は$\bB$の対角化行列であって、
      \begin{gather}
        \grL^{-1}\bB\grL
        = \bmat{
          -6 & 0 \\ 0 & 11
        }
      \end{gather}
      と対角化される。
      これを用いて$\bB^{n}$は
      \begin{align}
        \bB^{n}
        &= \grL \bmat{
          (-6)^{n} & 0 \\ 0 & 11^{n}
        }\grL^{-1} \\
        &= \frac{1}{17}\bmat{
          (-6)^{n} & 4\cdot 11^{n} \\
          -4\cdot(-6)^{n} & 11^{n}
        }\bmat{
          1 & -4 \\ 4 & 1
        } \\
        &= \frac{1}{17} \bmat{
          (-6)^{n} + 16\cdot 11^{n} &
          -4\cdot(-6)^{n} + 4\cdot 11^{n} \\
          -4\cdot(-6)^{n} + 4\cdot 11^{n} &
          16\cdot(-6)^{n} + 11^{n}
        }
      \end{align}
      よって、$\bB^{n}$の1行1列の成分$B^{n}_{\:11}$は
      \begin{gather}
        B^{n}_{\:11}
        = \frac{1}{17}\bigl(16\cdot 11^{n} + (-6)^{n}\bigr)
      \end{gather}
    \end{other*}
  \end{enumerate}
\end{ans*}
